\section{Related Work}


\iffalse
We first review existing work in scenario-based ADS testing and adversarial scenario generation, identifying the gap that motivates our approach (Sec.~\ref{sec:scenario_based_testing}--\ref{sec:scenario_generation}). We then introduce the simulation platforms and ADS used in our project (Sec.~\ref{sec:platform_ads}), followed by the two external frameworks --- the NHTSA pre-crash scenario typology and the Euro NCAP assisted driving protocol --- that ground our methodology in real-world crash statistics and regulatory kinematic standards (Sec.~\ref{sec:NHTSA}--\ref{sec:euro_ncap}).

\subsection{Scenario-Based Testing of ADS}\label{sec:scenario_based_testing}

Simulation-based testing has become the dominant paradigm for ADS evaluation. Kalra and Paddock~\cite{kalra_driving_2016} demonstrated that proving the statistical safety of an autonomous vehicle would require over $10^9$ miles of on-road driving --- an impractical requirement that has driven the community towards virtual, scenario-based evaluation. High-fidelity simulators such as CARLA~\cite{dosovitskiy_carla_2017} and LGSVL~\cite{rong_lgsvl_2020} provide controllable virtual environments with realistic sensor models, physics, and rendering, but their effectiveness depends critically on the quality, diversity, and realism of the test scenarios they execute.

Several naturalistic driving datasets have been developed to populate simulation testing pipelines. The inD dataset~\cite{bock_ind_2020} provides drone-recorded vehicle and pedestrian trajectories at German intersections, while highD~\cite{krajewski_highd_2018} captures highway traffic and rounD~\cite{krajewski_round_2020} records roundabout interactions. The Waymo Open Dataset~\cite{sun_scalability_2020} and nuScenes~\cite{caesar_nuscenes_2020} supply large-scale multi-sensor data from urban driving. The CARLA Leaderboard~\cite{noauthor_carla_nodate} defines ten traffic scenarios derived from the NHTSA pre-crash typology, providing a standardized benchmark for end-to-end ADS evaluation. The SCTrans dataset~\cite{dai_sctrans_2024}, which we use as the basis for our seed scenario identification, constructs over 1,900 simulation scenarios by transforming naturalistic traffic recordings into OpenSCENARIO files compatible with CARLA and LGSVL. SCTrans demonstrated that its scenarios can boost the collision-finding capability of existing ADS fuzzers by approximately seven times. However, a fundamental limitation shared by all these datasets is that they model \textit{well-behaved} traffic participants whose actions conform to expected driving norms, leaving adversarial edge cases underrepresented.

Scenario generation methods span three broad categories~\cite{wang_survey_2024, riedmaier_survey_2020, tang_survey_2023}. \textit{Rule-based methods} rely on expert knowledge, traffic regulations, and ontologies to derive test scenarios. Deng et al.~\cite{deng_target_2025} (TARGET) use LLM-guided knowledge extraction from traffic rules to generate test scenarios and evaluate seven ADS across CARLA, LGSVL, and MetaDrive. Birkemeyer et al.~\cite{birkemeyer_feature-interaction_2022} apply combinatorial testing with feature-interaction sampling for ADAS testing. While systematic, these approaches are constrained to scenarios that conform to predefined rules and expert expectations.

\textbf{Data-driven approaches} extract or synthesize scenarios from naturalistic data. TrafficGen~\cite{feng_trafficgen_2023} learns to generate diverse and realistic traffic scenarios, while TrafficComposer~\cite{tu_multi-modal_2025} combines natural language descriptions with traffic scene images for multi-modal scenario construction. Jia et al.~\cite{ji_effuzz_2023} use conditional generative adversarial imitation learning to produce different types of traffic scenarios. These methods are effective at reproducing realistic traffic distributions but inherit the bias of their training data --- since safety-critical events are extremely rare in naturalistic driving, data-driven methods tend to generate predominantly safe scenarios.

\textbf{Search-based methods} formulate generation as an optimization problem. AV-FUZZER~\cite{li_av-fuzzer_2020} uses genetic algorithms to search for NPC manoeuvres that lead to ADS safety violations, employing a fitness function based on safety potential. DriveFuzz~\cite{kim_drivefuzz_2022} extends fuzzing to the full ADS stack by mutating maps, actors, weather, and road conditions, discovering 30 bugs in Autoware and CARLA Behavior Agent. TM-fuzzer~\cite{lin_tm-fuzzer_2024} improves upon DriveFuzz by strategically managing NPC vehicle placement through a traffic management engine, achieving superior violation discovery. OSG~\cite{yan_-demand_2025} generates scenarios across varying risk levels using a Risk Intensity Regulator and speciation-based particle swarm optimization, providing a comprehensive ADS scoring system. EF-Fuzz~\cite{ji_effuzz_2023} introduces risk-guided seed selection, achieving a 37.9\% efficiency increase over DriveFuzz. ISS-Scenarios~\cite{li_iss-scenario_2024} provides a parameterized scenario library for batch testing in CARLA with both random sampling and genetic-algorithm-based search.

While effective at finding collision-inducing \textbf{configurations}, these search-based approaches treat external vehicle behaviour as a secondary search dimension --- mutating positions, speeds, or spawn locations rather than systematically modelling specific types of adversarial driving behaviour. The generated scenarios may expose bugs but often produce unrealistic NPC actions that do not correspond to plausible real-world driver behaviour.

\subsection{Adversarial and Safety-Critical Scenario Generation}\label{sec:scenario_generation}

A growing body of work specifically targets safety-critical scenarios through adversarial methods, which can be organized into three categories.
All three categories have the follwing key limitations: adversarial scenarios are very short in duration restricted to 5 seconds or less, narrow in scope focussingg on simple vehicle interactions in a very simple setting like a junction or single lane and finally lacking connection with realistic vehicle interactions that lead to collisions and other forms of undesirable behaviour. 

\textbf{Reinforcement learning approaches} train adversarial agents to challenge the ego vehicle. Learning-to-Collide~\cite{ding_learning_2020} represents traffic scenarios as autoregressive building blocks and trains a generative agent via policy gradient methods to search for risky scenario parameters. Feng et al.~\cite{feng_dense_2023} develop a dense deep reinforcement learning (D2RL) approach that edits Markov decision processes to densify safety-critical information, enabling naturalistic and adversarial environment generation. The CRASH framework~\cite{kulkarni_crash_2024} uses adversarial deep RL to control NPC agents, inducing collisions with the ego vehicle and iteratively refining the motion planner through safety hardening. AuthSim~\cite{yang_authsim_2025} proposes a three-layer relative safety region model combined with RL to ensure that generated scenarios are \textit{authentic} --- the ego vehicle bears primary collision responsibility rather than NPCs engaging in physically implausible attacks.

While RL-based approaches can discover effective adversarial strategies, they face several limitations: (1) they typically operate in narrow scenario classes (e.g., highway merging or single-lane following) rather than covering a comprehensive pre-crash typology; (2) they require extensive training time and computational resources; (3) the learned adversarial policies may produce behaviours that, while collision-inducing, do not correspond to recognizable real-world driving errors; and (4) they generally do not provide the kinematic interpretability needed to classify discovered vulnerabilities against established crash typologies.

\textbf{Trajectory perturbation approaches} directly modify actor trajectories. AdvSim~\cite{wang_advsim_2023} perturbs trajectories in a physically plausible manner and updates LiDAR sensor data to match the perturbed world, generating safety-critical scenarios for the full LiDAR-based autonomy stack. STRIVE~\cite{rempe_generating_2022} optimizes adversarial trajectories in the latent space of a VAE model to generate accident-prone driving scenarios. DiffScene~\cite{xu_diffscene_2025} applies diffusion models for safety-critical scenario generation, while CausalAF~\cite{ding_causalaf_2023} uses causal autoregressive flows. Game-theoretic methods~\cite{lin_scenario-based_2025} model adversarial vehicle interactions using kinematic constraints and reachable-set analysis. These methods can produce diverse scenarios but are typically limited to simplified driving environments, pairwise vehicle interactions, or lack systematic coverage of the full spectrum of pre-crash typologies.

\textbf{LLM and foundation model approaches} represent the most recent trend. ChatScene~\cite{zhang_chatscene_2024} uses knowledge-enabled LLMs for safety-critical scenario generation, outperforming prior methods across multiple ego vehicles and base scenarios. ScenGE~\cite{liu_adversarial_2025} combines LLM-based meta-scenario generation with retrieval-augmented grounding in traffic regulations and pre-crash scenario knowledge, using a structured knowledge base to ensure both plausibility and criticality. The retrieval-augmented framework by Nie et al.~\cite{mei_seeking_2025} employs an LLM-based behaviour analyser to infer the most dangerous intent of background vehicles and synthesize adversarial trajectories. SafeBench~\cite{xu_safebench_2022} provides a benchmarking platform integrating multiple generation algorithms (Learning-to-Collide, AdvSim, Adversarial Trajectory Optimization) with multiple ADS, though its predefined scenario types do not systematically cover the full NHTSA pre-crash typology.



\fi

A growing body of work targets safety-critical scenarios through adversarial methods, broadly spanning three categories. Across all categories, however, existing approaches share key limitations: scenarios are short in duration (typically $\leq5$ seconds), narrow in scope (e.g., simple junctions or single-lane interactions), and weakly grounded in realistic, complex vehicle behaviours that lead to real-world collisions.

Reinforcement learning approaches train adversarial agents to challenge the ego vehicle (e.g., Learning-to-Collide~\cite{ding_learning_2020}, D2RL~\cite{feng_dense_2023}, CRASH~\cite{kulkarni_crash_2024}, AuthSim~\cite{yang_authsim_2025}). While effective at inducing failures, they are typically confined to narrow scenario classes, require significant training cost, and often produce behaviours that lack realism or kinematic interpretability, limiting alignment with established crash typologies.

Trajectory perturbation approaches directly modify actor trajectories (e.g., AdvSim~\cite{wang_advsim_2023}, STRIVE~\cite{rempe_generating_2022}, DiffScene~\cite{xu_diffscene_2025}, CausalAF~\cite{ding_causalaf_2023}, game-theoretic methods~\cite{lin_scenario-based_2025}). Although they enable diverse scenario generation, they are generally restricted to simplified environments or pairwise interactions and do not provide systematic coverage of pre-crash scenarios.

LLM and foundation model approaches (e.g., ChatScene~\cite{zhang_chatscene_2024}, ScenGE~\cite{liu_adversarial_2025}, retrieval-augmented methods~\cite{mei_seeking_2025} , SafeBench~\cite{xu_safebench_2022}) improve diversity and plausibility through knowledge grounding. However, they still inherit limitations in scenario duration, interaction complexity, and systematic coverage of real-world pre-crash typologies.